% GENERATED, edit the French unit or the translation table
% GENERATED from content/data/enrichment.json by mxbuild.enrichment; do not edit.
\mxunit{enrichment-euler-trails-notes}{Routes without lifting the pencil}
\mxsubtitle{High school · graphs and logic \textperiodcentered{} Lesson and method \textperiodcentered{} 06}

\begin{prerequisites}
Vertices, edges and counting.
\end{prerequisites}

\begin{keypoint}[Ideas and method]
In an undirected graph, a vertex’s degree counts incident edge ends. A loop counts twice. The sum of the degrees is twice the number of edges.

An Euler trail uses each edge exactly once. At an intermediate vertex, each arrival must pair with a departure. Only the starting and ending vertices can therefore have odd degree.

If all vertices of nonzero degree are connected, an Euler trail exists when there are zero or two odd-degree vertices. With zero it ends where it starts; with two it runs between them. A trail can be built by joining cycles.
\end{keypoint}

\begin{example}
In a square ABCD with diagonal AC, degrees are $3,2,3,2$. A trail exists from A to C: A--B--C--D--A--C. It uses all five edges exactly once.
\end{example}

\begin{figurebox}
\begin{tikzpicture}[scale=1.1]\draw[mxmint,thick](0,0)--(2,0)--(2,2)--(0,2)--cycle;\draw[mxink,thick](0,0)--(2,2);\foreach \x/\y/\z in {0/0/A,2/0/B,2/2/C,0/2/D}{\fill[mxmint](\x,\y)circle(2pt)node[above right]{$\z$};}\end{tikzpicture}
\end{figurebox}

\begin{activity}[Explore]
Draw the square, then add the other diagonal. Recount degrees before trying a trail.
\end{activity}

